\documentclass[10pt,a4paper]{amsart}
\usepackage[margin=19mm]{geometry}
\usepackage{amsmath,amssymb,mathtools}
\usepackage[scr=rsfs,cal=euler]{mathalfa}
\usepackage{xfrac}
\usepackage[unicode,hidelinks,bookmarks=false]{hyperref}
\newcommand{\ie}{i.e.\ }
\newcommand{\cf}{cf.\ }
\newcommand{\resp}{respectively\ }
\newcommand{\Subsection}{\subsection}
\providecommand{\nobreakdash}{\nobreak\hbox{-}}
\setlength{\emergencystretch}{2em}
\multlinegap0pt
\usepackage{xcolor}
\usepackage{tikz}
\usepackage{float}
\usepackage{amssymb}
\usepackage{braket}
\usepackage{cancel}
\newtheorem{theorem}{Theorem}
\newtheorem{lemma}{Lemma}
\newtheorem{corollary}{Corollary}
\title{Topology of singular foliations of closed 1-forms on orbifolds}
\author{Daniel L\'opez-Garcia {\tiny and} Fabricio Valencia}
\date{Source: arXiv:2507.09644v2 (CC BY 4.0)}
\begin{document}
\maketitle

\noindent\emph{Source and licence.} Topology of singular foliations of closed 1-forms on orbifolds. Version arXiv:2507.09644v2 (CC BY 4.0), distributed by arXiv under the Creative Commons Attribution 4.0 International licence, http://creativecommons.org/licenses/by/4.0/, as recorded in the arXiv licence metadata for this exact version and shown on the abstract page https://arxiv.org/abs/2507.09644v2. Full licence notice: \texttt{licenses/arxiv-2507.09644-CC-BY-4.0.txt}.\par\smallskip

\noindent\textit{Editorial scope.} Counted unit: the article's corollary cor1 (source0.tex lines 423--426). The article's complete original proof of it is delivered with it (source0.tex lines 428--431, one \texttt{proof} environment of 4 lines). No mathematics is written, repaired or re-derived: the statement and the proof are the article's own lines, in the article's own notation, and no retained line is substituted or re-flowed.  Retained uncounted as the source-local context the proof needs: lines 156--173; lines 331--334.  Nothing was omitted from any retained span, so the package carries no omission record.  No retained line contains a citation, so the wrapper carries no bibliography and no bibliography entry.  No retained line contains a figure environment, an image-inclusion command or drawing code, so no image is delivered and the package declares no asset dependency.  Editorial formatting only, in addition: an AMS \texttt{amsart} preamble that restates the article's own declarations and macros used by the retained lines, namely the environments \texttt{theorem}, \texttt{lemma}, \texttt{corollary} on separate counters, together with the standard packages those lines need.  The retained environments print in this document as Theorem 1, Lemma 1, Corollary 1 in order of appearance, whereas the article prints the same items with the numbering of its own full document.  The complete native source of the article as distributed in the arXiv e-print for this exact version is bundled unchanged as \texttt{sources/181495/source0.tex}.
\begin{theorem}\label{thm1}
Let $X$ be a compact orbifold of dimension $n$ and let $\overline{\omega}$ denote a closed 1-form of Morse type on $X$ representing a cohomology class $\xi\in H^1(X)$.
\begin{enumerate}
\item All leaves of the singular foliation $\tilde{\mathcal{F}}_\omega$ are
compact if and only if the homomorphism of periods $\Pi_1^{\textnormal{orb}}(X)\to (\mathbb{R},+)$ induced by $\xi$ can be factorized 
as 
$$\Pi_1^{\textnormal{orb}}(X)\to F\to (\mathbb{R},+),$$
through a free group $F$. 
\item $X$ is the union of two compact $n$-dimensional suborbifolds $X_c$ and $X_\infty$ with a common, possibly singular, boundary satisfying:
\begin{enumerate}
	\item $\textnormal{int}(X_c)$ is the union of all of the compact leaves (nonsingular and singular) plus some compact singular leaf components of non-compact singular leaves,
	\item $\textnormal{int}(X_\infty)$ is the union of all noncompact nonsingular leaves, all noncompact singular leaf components and some compact singular leaf components of noncompact singular leaves,
	\item $\partial  X_c=\partial X_\infty=X_c\cap X_\infty$ is the union of finitely many compact singular leaf components of noncompact singular leaves. It is an orbifold except at finitely many points which are zeros of $\overline{\omega}$,
	\item if $X_\infty\neq \emptyset$ then the rank of the cohomology class $\xi_\infty=\xi|_{X_\infty}\in H^1(X_\infty)$ is greater than $1$, and
	\item if $\overline{\omega}$ is generic then the boundary $\partial  X_c=\partial X_\infty=X_c\cap X_\infty$ is the union of all compact singular leaf components of non-compact singular leaves. It is a closed $(n-1)$-dimensional topological orbifold.
\end{enumerate}
\end{enumerate}
\end{theorem}

\begin{lemma}\label{w-graph}
The singular leaves are isolated (i.e. each singular leaf has
a neighborhood, free of other singular leaves) and the quotient space $\overline{\Gamma}_\omega=X/\sim$ is a graph.
\end{lemma}

\begin{corollary}\label{cor1}
Let $\overline{\omega}$ be a closed 1-form of Morse type on a compact and connected orbifold $X$ with no zeros of indices $1$ and $n-1$. Then, either 
any leaf of the singular foliation $\tilde{\mathcal{F}}_\omega$ is compact or any leaf is dense in $X$. If $\xi=[\omega]\in H^1(X)$ then the first case happens if and only if $\textnormal{rk}(\xi)\leq 1$ and the second case happens if and only if $\textnormal{rk}(\xi)>1$.
 \end{corollary}

\begin{proof}
If $\overline{\omega}$ has no zeros of indices $1$ and $n-1$ then the boundary $X_c\cap X_\infty$ is empty, so that either $X_c=\emptyset$ (all leaves are noncompact) or else $X_\infty=\emptyset$ (all leaves are compact). We only have to prove that $\textnormal{rk}(\xi)\leq 1$ if all leaves are compact, see item (d) in Theorem \ref{thm1}. By Lemma \ref{w-graph} we obtain that $\overline{\Gamma}_\omega$ is homemorphic to a circle or a closed interval, as all of its vertices are bivalent or 
univalent in this case. Hence, the results follows by applying Theorem \ref{thm1}.
\end{proof}

\end{document}
